
Research pipelines fail silently when semantic contradictions pass syntactic checks. Existing validation tools catch schema violations but miss reasoning failures because they treat execution traces as opaque logs rather than semantic artifacts. This work demonstrates that Model Context Protocol (MCP) tool-calling traces encode researcher reasoning in natural language, enabling automated semantic extraction. A three-layer framework (syntactic validation, semantic extraction, constraint inference) was evaluated on 20 research pipeline executions containing 596 tool calls. Layers 1-2 achieved 97.48\% trace completeness, 97.48\% natural language presence in both queries and results, and 82.7\% extraction recall with 86.3\% precision (Cohen's $\kappa$=0.716) using zero-shot LLM prompting without manual annotation. Layer 3 constraint inference via semantic similarity matching failed to detect contradictions (0\% recall across 1,200 assumption-claim pairs), revealing that sentence embeddings optimize for topic relatedness rather than logical contradiction. The validated two-layer framework establishes zero-annotation feasibility for syntactic validation and semantic extraction in research pipelines, while the negative result on constraint inference identifies a boundary condition and points toward entailment models or LLM-based reasoning as necessary alternatives.
